%% kapa: exam=1 round=27 year=2016 period=1 form=A subject=economics subject_ko=경제학원론 seq=7 number=47
%% source: Q-Net 감정평가사 기출문제 (cst00309) artlSeq=5169402 file=extracted/1008/5169402/제27회 감정평가사 1차 문제/제27회 감정평가사 1교시문제(A형).pdf
%% conversion: kordoc 4.24.51 -> tools/gen1.py (auto); review=proofread-sonnet; flags=none
\begin{question}{47}
소비자 이론에 관한 설명으로 옳은 것은? (단, 소비자는 X재와 Y재만 소비한다.)
\begin{choices}
  \item 소비자의 효용함수가 $U=2XY$일 때, 한계대체율은 체감한다.
  \item 소비자의 효용함수가 $U=\sqrt{XY}$ 일 때, X재의 한계효용은 체증한다.
  \item 소비자의 효용함수가 $U=\min(X,Y)$일 때, 수요의 교차탄력성은 0이다.
  \item 소비자의 효용함수가 $U=\min(X,Y)$일 때, 소득소비곡선의 기울기는 음(-)이다.
  \item 소비자의 효용함수가 $U=X+Y$일 때, X재의 가격이 Y재의 가격보다 크더라도 X재와 Y재를 동일 비율로 소비한다.
\end{choices}
\end{question}
\answer{1}
